% !TeX root = ../../main.tex
\section{\bilingual{Related Work}{研究现状}}

    \bilingual{
        The study of extremal behavior in BRW has evolved from first-order asymptotic speed analysis to fine fluctuation analysis. Since the 1970s, Hammersley~\mycite{hammersley1974postulates}, Kingman~\mycite{kingman1975first}, and Biggins~\mycite{biggins1976first} established the first-order asymptotic speed of the extremal position in BRW. Subsequently, Bramson~\mycite{Bramson1978MaximalDO} discovered in his classic work on BBM that the front of the maximum position does not advance uniformly but has a famous second-order logarithmic correction term $-\frac 32\log n$, indicating that the extremal behavior of BRW has a more complex asymptotic structure than independent random walks.
    }{
        BRW 极值理论的研究经历了从一阶渐近速度到精细涨落分析的发展过程. ~自 20 世纪 70 年代起,~ Hammersley~\mycite{hammersley1974postulates},~ Kingman~\mycite{kingman1975first} 与 Biggins~\mycite{biggins1976first} 确立了 BRW 最值位置的一阶渐近速度. ~随后,~ Bramson~\mycite{Bramson1978MaximalDO} 在 BBM 的经典工作中发现,~ 最大位置前沿并非均匀推进,~ 而是存在一个著名的次阶对数修正项 $-\frac 32\log n$,~ 表明 BRW 的极值行为具有比独立随机游走更复杂的渐近结构.
    }

    \bilingual{
        The shift in research paradigm in this field largely relies on the innovation of probabilistic analysis tools, especially the maturation of the Lyons--Pemantle--Peres--Biggins spine decomposition theory~\mycite{lyons1995conceptual, biggins2004measure}. By constructing an infinite spine under change of measure, researchers have been able to transform the complex summation problem in tree space into the analysis of a single-path random walk (\nameref{lemma:many-to-one}), allowing for the use of large deviation and small deviation theories of random walks to precisely characterize the extremal properties of tree structures. This provides a core mathematical framework for quantitative analysis from macroscopic speed to microscopic structure.
    }{
        这一领域研究范式的转向,~ 很大程度上依赖于概率分析工具的革新,~ 特别是 Lyons--Pemantle--Peres--Biggins 脊柱分解 (spine decomposition) 理论\mycite{lyons1995conceptual, biggins2004measure} 的成熟. ~通过构造测度变换下的无限长脊柱,~ 研究者得以将复杂的树状空间求和问题转化为单路径随机游走的分析 (\nameref{lemma:many-to-one}),~ 这一方法允许利用随机游走的大偏差理论和小偏差理论来精确刻画树状结构的极值性质. ~这为从宏观速度转向微观结构的定量解析提供了核心数学框架.
    }

    \bilingual{
        Under this framework,~ the research focus after the 21st century has shifted to the weak convergence properties of extremal particles. A\"id\'ekon~\mycite{A_d_kon_2013} combined the derivative martingale technique with change of measure methods to prove that the extremal position has a non-degenerate weak limit after second-order correction. Subsequently, Madaule~\mycite{madaule2017convergence} and others further characterized the extremal process, proving its convergence to a decorated Poisson point process (DPPP), thus completing a comprehensive mathematical description of the microscopic structure of BRW extremes.
    }{
        这一框架下,~ 21 世纪后研究重点转向极值的弱收敛性质. ~A\"id\'ekon~\mycite{A_d_kon_2013} 结合了{\kaishu 导数鞅} (derivative martingale) 技术与测度变换的方法,~ 证明了极值位置在经过次阶修正后具有非退化的弱极限. ~此后,~ Madaule~\mycite{madaule2017convergence} 等人进一步刻画了极值点过程 (extremal process),~ 证明其收敛至一个{\kaishu 装饰Poisson点过程} (decorated Poisson point process, DPPP),~ 完成了对 BRW 极值微观结构的完整数学描述.
    }

    \bilingual{
        As for the $N$-BRW model with selection mechanism,~ the core issue is to verify a conjecture proposed by Brunet and Derrida et al.~\mycite{Brunet_2007} regarding the velocity correction term. The physical idea behind this conjecture is as follows: in the deterministic KPP equation, the wave speed stabilizes at $v_\infty$; however, in the finite-$N$ $N$-BRW, the cutoff of the left tail of the front acts as an effective noise source that persistently perturbs the front position. Due to the exponential amplification effect of branching, this microscopic cutoff leads to a macroscopic slowdown of the wave speed compared to $v_\infty$; moreover, the magnitude of this slowdown is logarithmically sensitive to $N$, asymptotically scaling as $(\log N)^{-2}$. The conjecture takes the specific form
    }{
        对于带选择机制的 $N$-BRW 模型,~ 核心课题是验证 Brunet 与 Derrida 等\mycite{Brunet_2007} 提出的关于速度修正项的猜想. 该猜想的物理思想如下：在确定性 KPP 方程中, 波速稳定于 $v_\infty$; 然而, 在有限 $N$ 的 $N$-BRW 中, 前沿左尾截断作为一种有效噪声源持续扰动前沿位置. 由于分支过程对前沿扰动具有指数级的放大效应, 这种微观截断导致宏观波速相比 $v_\infty$ 产生可观测的减慢; 并且,~ 减慢幅度对 $N$ 具有对数级的敏感性,~ 渐近标度为 $(\log N)^{-2}$. 该猜想的具体形式为
    }
    \begin{equation*}
        \frac{C}{\left(\log N+3\log\log N+o_N(\log\log N)\right)^2},
    \end{equation*}
    \bilingual{
        where $C$ is an explicit constant determined by the branching law. The first challenge is to rigorously establish the $(\log N)^{-2}$ scaling law of the velocity correction term with respect to the population size $N$. Bérard and Gouéré~\mycite{bérard2010brunetderridabehaviorbranchingselectionparticle} were the first to mathematically prove the correctness of this asymptotic order. The applicability of this result is limited to binary branching models where each generation produces two identically distributed offspring. Subsequently, Mallein~\mycite{Mallein_2017} further extended the study to the convergence of $v_N$ under general $\alpha$-stable displacement distributions, establishing its universal $(\log N)^{-2}$ scaling law. Secondly, the formula implies the existence of a secondary correction term, namely $(\log N)^{-3}\log\log N$.~ Although this fine term has a clear physical meaning (reflecting the logarithmic delay characteristic of the front particle {\kaishu cloud}), mathematically, except for a few specific soft selection models that can be solved exactly~\mycite{cortines2016nbranchingrandomwalkrandom}, proving this higher-order expansion term under the original hard selection $N$-BRW framework remains an open problem.
    }{
        其中$C$是由分支律决定的显式常数. ~首先面临的问题是严格确立速度修正项相对于种群规模 $N$ 的 $(\log N)^{-2}$ 标度律. Bérard 与 Gouéré~\mycite{bérard2010brunetderridabehaviorbranchingselectionparticle} 首次在数学上证明了该渐近阶数的正确性. ~该结果的适用范围限定于每代产生两个同分布子代的二分支模型. ~此后,~ Mallein~\mycite{Mallein_2017} 进一步将研究拓展到一般$\alpha$-稳定位移分布下 $v_N$ 的收敛性,~ 确立了其普遍的$(\log N)^{-2}$ 标度律. ~其次, ~该公式暗示了次级修正项的存在,~ 即 $(\log N)^{-3}\log\log N$.~ 虽然这一精细项在物理学中具有明确的含义 (反映了前沿粒子{\kaishu 云团} (cloud) 的对数延迟特征), 但在数学上,~ 除了少数可精确求解的特定软选择模型\mycite{cortines2016nbranchingrandomwalkrandom} 外,~ 在原始的硬选择 $N$-BRW 框架下证明该高阶展开项仍然是一个开放问题.
    }

    \bilingual{
        Another research direction for $N$-BRW is its genealogical structure traced back in time. In neutral evolutionary models without selection pressure, the {\itshape genealogy} of the population typically follows the {\itshape Kingman coalescent}; however, in $N$-BRW with strong selection pressure, many individuals in the current generation share an extremely outstanding {\itshape super-ancestor}, leading to multiple branches coalescing simultaneously when tracing back in time. This feature of {\itshape multiple collisions} corresponds to the {\itshape Bolthausen--Sznitman coalescent} (BS coalescent), which was first predicted by physicists~\mycite{Brunet_2007, Brunet_2012}. Mathematically, Berestycki et al.~\mycite{Berestycki_2013} rigorously proved this conclusion for $N$-BBM with absorbing boundaries. However, for $N$-BRW under hard selection mechanism, the asymptotic behavior of its genealogical structure and the rigorous proof of the coalescence time scale still face technical difficulties.
    }{
        $N$-BRW 的另一研究方向是其随时间回溯的系谱结构. ~在无选择压力的中性演化模型中,~ {\kaishu 系谱} (genealogy)通常遵循 {\kaishu Kingman 聚合} (Kingman coalescent);~ 但在具有强选择压力的$N$-BRW 中,~ 当前某一代的很多个体都共有一个极其出色的{\kaishu 超级祖先} (super-ancestor),~ 使得在回溯时间时,~ 多个分支会同时聚合,~ 这种具有{\kaishu 多体碰撞} (multiple collisions) 特征的 {\kaishu Bolthausen--Sznitman 聚合} (BS coalescent) 首先被物理学家预言\mycite{Brunet_2007, Brunet_2012}. ~在数学层面,~ Berestycki 等\mycite{Berestycki_2013} 为带吸收边界的$N$-分支Brownian运动严格证明了这一结论. ~然而,~ 对于硬选择机制下的 $N$-BRW,~ 其系谱结构的渐近行为以及合并时间尺度的严格证明仍面临技术困难.
    }

    \bilingual{
        Recently, the research focus on $N$-BRW has shifted from order estimation of macroscopic speed to fine characterization of frontier dynamic trajectories and edge fluctuations. A\"id\'ekon and Hu~\mycite{aïdékon2025averageweightpercolationcompletegraph} studied a class of $N$-branching random walks with a specific Poisson structure, establishing an $\mathrm L^2$ limit theorem for the position of the leftmost particle under asymmetric scaling of log-cubed time and single logarithmic space, proving that its position follows a linear trajectory during evolution, thus quantitatively revealing the influence of initial population size on the evolutionary path of the frontier, providing a new mathematical framework for understanding the stability of constrained branching systems.
    }{
        近期,~ 关于 $N$-BRW 的研究重心已从宏观速度的阶数估计转向对前沿动态轨迹及边缘涨落的精细刻画. ~A\"id\'ekon 与 Hu \mycite{aïdékon2025averageweightpercolationcompletegraph} 针对一类具有特定 Poisson 结构的 $N$-分支随机游走,~ 在三次对数时间与一次对数空间的非对称缩放尺度下,~ 确立了最左侧粒子位置的 $\mathrm L^2$极限定理,~ 证明其位置在演化过程中遵循线性轨迹,~ 从而定量揭示了初始种群规模对前沿演化路径的影响,~ 这为理解受限分支系统的稳定性提供了新的数学框架.
    }
